\cotefile{157-3}
\scanpage{312}
\mspage{262}

carré $W$-cocartésien, qui paraît \unsure{inépuisable} : Deux points de vue, et d'une simplicité parfaite.

Cela amène à revoir la notion en question dans $\mathrm{Cat}$, où elle avait l'air très \unsure{astucieuse}, profonde et un peu obscure, \unsure{un brin vertigineuse|vu du vertiguiste|nom de vertiguiste} et tout ça. C'est le moment de se rendre compte que ``ce n'est que ça''. Donc de prouver la

\uline{Proposition 20} Supposons $\mathcal M=\mathrm{Cat}$, $W\unsure{=}\text{un}$ localiseur fondamental satisfaisant $W\,(7\ \mathrm{bis})$.

Considérons un carré commutatif dans $\mathrm{Cat}$

\diagram{\deleted{\(Q=\left(\begin{tikzcd}[ampersand replacement=\&,column sep=small,row sep=small]
X_{\unsure{1}}\arrow[r,leftrightarrow] \& X_3\arrow[d,leftrightarrow]\\
S\arrow[u,"{\unsure{p_1|p'_1}}"]\arrow[r,leftrightarrow,"{\unsure{p'_2|p_2}}"'] \& X_2
\end{tikzcd}\right).\)}
\qquad \(Q=\left(\begin{tikzcd}[ampersand replacement=\&,column sep=small,row sep=small]
S\arrow[r,dash]\arrow[d] \& X_1\arrow[d]\\
X_2\arrow[r,dash] \& X_3
\end{tikzcd}\right).\)}

Pour qu'il soit $W$-cartésien au sens de la théorie \deleted{\unsure{générale}} développée ici (pour des catégories de modèles générales), il f. et s. que

\diagram{\(\varphi_Q:\left(\overbrace{\int\ \deleted{\substack{\unsure{X_1}\\\unsure{X_2}}}\begin{tikzcd}[ampersand replacement=\&,column sep=tiny,row sep=tiny]
\& X_1\arrow[dl]\\ S\&{}\\ \& X_2\arrow[ul]
\end{tikzcd}}^{\unsure{\mathscr X_Q|X_Q|\mathcal X_Q}}\right)\longrightarrow \deleted{\unsure{D}}X_3\)}

soit dans $W$. (NB L'intérêt de ce critère est qu'il est ``intrinsèque'' : ne faisant pas appel à une factorisation des $p_i$ en des \unsure{fl} en

\diagram{\(\deleted{\unsure{p_i}\longrightarrow\widetilde X_i}\quad\begin{tikzcd}[ampersand replacement=\&,column sep=large]
S\arrow[r,"{\in\mathrm{Cof}\deleted{\unsure{\cap W}}}"'] \& \widetilde X_i\arrow[r,"{\in W}"'] \& X_i
\end{tikzcd}.\unsure{\}}\)}

Cf. applications intéressantes dans section V 3.4. p. 146 ff \unsure{(th.|th.} 4 et corollaires.)
% archive sheet number: 312
